\section{模型评价与改进}

\subsection{模型优点}

\begin{enumerate}[label=\arabic*)]
  \item 采用实际处理量和库存守恒，严格区分“到达前空闲产能”与“到达后可用产能”；
  \item 班次起点、人数、低效率小时和月度出勤均由整数规划统一刻画，结果可解释且可验证；
  \item 问题三同时给出下界证明和具体排班矩阵，不仅报告人数，还提供可执行人员编号方案；
  \item 程序自动生成结果表、图和论文引用文件，减少手工抄写引起的数值不一致。
\end{enumerate}

\subsection{局限性与改进方向}

模型假设班次按整数小时开始且不跨日；若允许跨日，可将时间轴扩展为连续48小时并增加相邻日联动。模型还未显式考虑设备数量、工人技能、交接损失和随机缺勤。实际应用中可引入鲁棒优化或情景规划，将进货预测误差和缺勤率纳入安全容量；也可增加不同班次成本、夜班补贴和公平性目标，形成多目标人员排班模型。

\section{结论}

本文针对物流分拣中心三问建立了由日内货物流模型和月度人员模型组成的混合整数规划体系。问题一最少需要10791人日，单日峰值537人；问题二在提前截止与低效率工时条件下需要11951人日，峰值581人；问题三通过理论下界与完整可行排班共同证明最少招工581人。处理效率下降5\%时问题二总人日上升至12576，说明人员需求对作业效率近似反比例但受离散班次约束放大。最终方案满足题目全部约束，并提供可复现代码、逐日班次、581名工人的出勤表和班次分配表。
